% GENERATED FILE. Edit canonical lessons, then run npm run generate:books.
% canonical-source-hash: fnv1a64:3bf1ae194c04f8d2
% canonical-lessons: PA-C299-kandha, PA-C299-khal, PA-C299-lahir, PA-C299-khara, PA-C299-sohna

\chapter{By the Water}
\label{ch:pa-a2-by-the-water}
\hlchaptermodality{299}

\begin{chapteropening}
\textbf{By the end of this chapter:} \emph{I can describe a lake and a riverbank.}

Complete the last lesson of chapter 299: I can describe a lake and a riverbank.
\end{chapteropening}

\section[\texorpdfstring{\pa{ਕੰਢਾ} (kaṇḍhā)}{ਕੰਢਾ (kaṇḍhā)}]{\pa{ਕੰਢਾ} (kaṇḍhā) — a riverbank; a shore}

\label{lesson:PA-C299-kandha}



\begin{quote}
Before the new one: say the Punjabi for dense, then the Punjabi for to wilt.
\end{quote}

\subsection*{You'll want to know: \pa{ਕੰਢਾ}}
\textbf{\pa{ਕੰਢਾ}} — \emph{kaṇḍhā} — \textquotedblleft{}a riverbank; a shore\textquotedblright{}.

\begin{quote}
\pa{ਅਸੀਂ} \pa{ਦਰਿਆ} \pa{ਦੇ} \pa{ਕੰਢੇ} \pa{ਉੱਤੇ} \pa{ਬੈਠਦੇ} \pa{ਹਾਂ।}
\end{quote}

\begin{quote}
\emph{asīṁ dariā de kaṇḍhe utte baiṭhde hāṁ.}
\end{quote}

\begin{quote}
We sit on the bank of the river.
\end{quote}

A masculine noun.

\begin{cousinweb}
\textbf{\pa{ਕੰਢਾ}} (\emph{kaṇḍhā}) is inherited from Sanskrit \emph{kaṇṭha}, \textquotedblleft{}throat, neck\textquotedblright{}, also \textquotedblleft{}border\textquotedblright{}; CDIAL lists Punjabi \emph{kaṇḍhā}, \textquotedblleft{}bank, shore\textquotedblright{}, under it (CDIAL 2680).
\end{cousinweb}

\begin{grammarlens}[title={What you've built}]
One more word for this chapter.
\end{grammarlens}

\subsection*{Your turn}
\begin{itemize}
\raggedright
  \item \emph{Say it:} \emph{kaṇḍhā}
  \item \emph{Say it:} \emph{kaṇḍhā}, once more
  \item \emph{Say it:} \emph{asīṁ dariā de kaṇḍhe utte baiṭhde hāṁ.} — We sit on the bank of the river.
  \item \emph{Recall:} say the Punjabi for dense, then the Punjabi for to wilt, then say \emph{kaṇḍhā} again
\end{itemize}

\begin{tcolorbox}[breakable,colback=teal!4,colframe=teal!35!black,title={Before you move on}]
What does \pa{ਕੰਢਾ} mean? (\textquotedblleft{}A riverbank; a shore\textquotedblright{}.) Say it once more.
\end{tcolorbox}
\section[\texorpdfstring{\pa{ਖਾਲ} (khāl)}{ਖਾਲ (khāl)}]{\pa{ਖਾਲ} (khāl) — a watercourse; a field channel}

\label{lesson:PA-C299-khal}



\begin{quote}
Before the new one: say the Punjabi for to wilt, then the Punjabi for a riverbank.
\end{quote}

\subsection*{You'll want to know: \pa{ਖਾਲ}}
\textbf{\pa{ਖਾਲ}} — \emph{khāl} — \textquotedblleft{}a watercourse; a field channel\textquotedblright{}.

\begin{quote}
\pa{ਖਾਲ} \pa{ਵਿੱਚ} \pa{ਪਾਣੀ} \pa{ਖੇਤ} \pa{ਤੱਕ} \pa{ਜਾਂਦਾ} \pa{ਹੈ।}
\end{quote}

\begin{quote}
\emph{khāl vich pāṇī khet tak jāndā hai.}
\end{quote}

\begin{quote}
In the channel the water goes to the field.
\end{quote}

A masculine noun.

\begin{cousinweb}
\textbf{\pa{ਖਾਲ}} (\emph{khāl}) is inherited from Sanskrit \emph{khalla}, \textquotedblleft{}trench, creek, canal\textquotedblright{}; CDIAL lists Punjabi \emph{khāl}, \textquotedblleft{}hollow, gutter, inlet\textquotedblright{}, under it (CDIAL 3849).
\end{cousinweb}

\begin{grammarlens}[title={What you've built}]
One more word for this chapter.
\end{grammarlens}

\subsection*{Your turn}
\begin{itemize}
\raggedright
  \item \emph{Say it:} \emph{khāl}
  \item \emph{Say it:} \emph{khāl}, once more
  \item \emph{Say it:} \emph{khāl vich pāṇī khet tak jāndā hai.} — In the channel the water goes to the field.
  \item \emph{Recall:} say the Punjabi for to wilt, then the Punjabi for a riverbank, then say \emph{khāl} again
  \item \emph{Write it:} copy \textbf{\pa{ਖਾਲ}} once, checking the whole word
\end{itemize}

\begin{tcolorbox}[breakable,colback=teal!4,colframe=teal!35!black,title={Before you move on}]
What does \pa{ਖਾਲ} mean? (\textquotedblleft{}A watercourse; a field channel\textquotedblright{}.) Say it once more.
\end{tcolorbox}
\section[\texorpdfstring{\pa{ਲਹਿਰ} (lahir)}{ਲਹਿਰ (lahir)}]{\pa{ਲਹਿਰ} (lahir) — a wave}

\label{lesson:PA-C299-lahir}



\begin{quote}
Before the new one: say the Punjabi for a riverbank, then the Punjabi for a field channel.
\end{quote}

\subsection*{You'll want to know: \pa{ਲਹਿਰ}}
\textbf{\pa{ਲਹਿਰ}} — \emph{lahir} — \textquotedblleft{}a wave\textquotedblright{}.

\begin{quote}
\pa{ਝੀਲ} \pa{ਵਿੱਚ} \pa{ਛੋਟੀਆਂ} \pa{ਲਹਿਰਾਂ} \pa{ਹਨ।}
\end{quote}

\begin{quote}
\emph{jhīl vich chhoṭīāṁ lahirāṁ han.}
\end{quote}

\begin{quote}
There are small waves on the lake.
\end{quote}

A feminine noun.

\begin{cousinweb}
\textbf{\pa{ਲਹਿਰ}} (\emph{lahir}) is inherited from Sanskrit \emph{laharī}, \textquotedblleft{}large wave\textquotedblright{}; CDIAL lists Punjabi \emph{lahir}, \textquotedblleft{}wave\textquotedblright{}, under it (CDIAL 10999).
\end{cousinweb}

\begin{grammarlens}[title={What you've built}]
One more word for this chapter.
\end{grammarlens}

\subsection*{Your turn}
\begin{itemize}
\raggedright
  \item \emph{Say it:} \emph{lahir}
  \item \emph{Say it:} \emph{lahir}, once more
  \item \emph{Say it:} \emph{jhīl vich chhoṭīāṁ lahirāṁ han.} — There are small waves on the lake.
  \item \emph{Recall:} say the Punjabi for a riverbank, then the Punjabi for a field channel, then say \emph{lahir} again
  \item \emph{Write it:} copy \textbf{\pa{ਲਹਿਰ}} once, checking the whole word
\end{itemize}

\begin{tcolorbox}[breakable,colback=teal!4,colframe=teal!35!black,title={Before you move on}]
What does \pa{ਲਹਿਰ} mean? (\textquotedblleft{}A wave\textquotedblright{}.) Say it once more.
\end{tcolorbox}
\section[\texorpdfstring{\pa{ਖਾਰਾ} (khārā)}{ਖਾਰਾ (khārā)}]{\pa{ਖਾਰਾ} (khārā) — brackish; salty}

\label{lesson:PA-C299-khara}



\begin{quote}
Before the new one: say the Punjabi for a field channel, then the Punjabi for a wave.
\end{quote}

\subsection*{You'll want to know: \pa{ਖਾਰਾ}}
\textbf{\pa{ਖਾਰਾ}} — \emph{khārā} — \textquotedblleft{}brackish; salty\textquotedblright{}.

\begin{quote}
\pa{ਸਮੁੰਦਰ} \pa{ਦਾ} \pa{ਪਾਣੀ} \pa{ਖਾਰਾ} \pa{ਹੁੰਦਾ} \pa{ਹੈ।}
\end{quote}

\begin{quote}
\emph{samundar dā pāṇī khārā hundā hai.}
\end{quote}

\begin{quote}
Sea water is salty.
\end{quote}

\begin{cousinweb}
\textbf{\pa{ਖਾਰਾ}} (\emph{khārā}) is inherited from Sanskrit \emph{kṣāra}, \textquotedblleft{}corrosive\textquotedblright{}; CDIAL lists Punjabi \emph{khārā}, \textquotedblleft{}brackish\textquotedblright{}, under it (CDIAL 3674).
\end{cousinweb}

\begin{grammarlens}[title={What you've built}]
One more word for this chapter.
\end{grammarlens}

\subsection*{Your turn}
\begin{itemize}
\raggedright
  \item \emph{Say it:} \emph{khārā}
  \item \emph{Say it:} \emph{khārā}, once more
  \item \emph{Say it:} \emph{samundar dā pāṇī khārā hundā hai.} — Sea water is salty.
  \item \emph{Recall:} say the Punjabi for a field channel, then the Punjabi for a wave, then say \emph{khārā} again
  \item \emph{Write it:} copy \textbf{\pa{ਖਾਰਾ}} once, checking the whole word
\end{itemize}

\begin{tcolorbox}[breakable,colback=teal!4,colframe=teal!35!black,title={Before you move on}]
What does \pa{ਖਾਰਾ} mean? (\textquotedblleft{}Brackish; salty\textquotedblright{}.) Say it once more.
\end{tcolorbox}
\section[\texorpdfstring{\pa{ਸੋਹਣਾ} (sohṇā)}{ਸੋਹਣਾ (sohṇā)}]{\pa{ਸੋਹਣਾ} (sohṇā) — beautiful; lovely}

\label{lesson:PA-C299-sohna}



\begin{quote}
Before the new one: say the Punjabi for a wave, then the Punjabi for brackish.
\end{quote}

\subsection*{You'll want to know: \pa{ਸੋਹਣਾ}}
\textbf{\pa{ਸੋਹਣਾ}} — \emph{sohṇā} — \textquotedblleft{}beautiful; lovely\textquotedblright{}.

\begin{quote}
\pa{ਸ਼ਾਮ} \pa{ਨੂੰ} \pa{ਝੀਲ} \pa{ਬਹੁਤ} \pa{ਸੋਹਣੀ} \pa{ਹੁੰਦੀ} \pa{ਹੈ।}
\end{quote}

\begin{quote}
\emph{shām nū̃ jhīl bahut sohṇī hundī hai.}
\end{quote}

\begin{quote}
In the evening the lake is very beautiful.
\end{quote}

\begin{cousinweb}
\textbf{\pa{ਸੋਹਣਾ}} (\emph{sohṇā}) is inherited from Sanskrit \emph{śōbhana}, \textquotedblleft{}brilliant, beautiful\textquotedblright{}; CDIAL lists Punjabi \emph{sohaṇā}, \textquotedblleft{}beautiful\textquotedblright{}, under it (CDIAL 12637).
\end{cousinweb}

\begin{grammarlens}[title={What you've built}]
One more word for this chapter.
\end{grammarlens}

\subsection*{Your turn}
\begin{itemize}
\raggedright
  \item \emph{Say it:} \emph{sohṇā}
  \item \emph{Say it:} \emph{sohṇā}, once more
  \item \emph{Say it:} \emph{shām nū̃ jhīl bahut sohṇī hundī hai.} — In the evening the lake is very beautiful.
  \item \emph{Recall:} say the Punjabi for a wave, then the Punjabi for brackish, then say \emph{sohṇā} again
  \item \emph{Write it:} copy \textbf{\pa{ਸੋਹਣਾ}} once, checking the whole word
\end{itemize}

\begin{tcolorbox}[breakable,colback=teal!4,colframe=teal!35!black,title={Before you move on}]
What does \pa{ਸੋਹਣਾ} mean? (\textquotedblleft{}Beautiful; lovely\textquotedblright{}.) Say it once more.
\end{tcolorbox}
